\section{The contract, and what it rules out}\label{sec:contract}

Every work rate in this paper is an optimal asymptotic rate at vanishing error over a complete interaction, with hardware allowed to grow. Such a statement is only as good as the contract behind it, so we state the contract first, as a threat model: what the device is allowed, what the observer is allowed, and which loopholes the rules close. Appendix~\ref{app:glossary} defines each term in one line.

\paragraph{The device.} A device consists of a controller $C$, a mechanical battery $W$ and Gibbs reservoirs $R$ at inverse temperature $\beta$, all declared in advance and initially independent. The controller holds all hardware: memory, clocks, catalysts, routing registers, records and any workspace not charged to the battery. Every operation is a unitary that conserves the total energy of the systems it acts on. There is no external source of coherence and no Gibbs-preserving map supplied as a black box. Measurements and records are part of the dilation and are counted. Work is the battery's mean loss of energy, and the battery's entropy may change only by $o(n)$, so that it cannot hide a second bath. Work is \emph{deterministic} when the battery starts in an energy eigenstate, ends in a fixed level with probability tending to one, and its mean absolute energy deviation from that level is $o(n)$. The battery is \emph{stationary} if its initial state commutes with its Hamiltonian; a battery in a broad superposition of energy levels would itself be a clock, which the device could spend without returning it.

\paragraph{The observer.} The observer supplies the carriers and may choose them adaptively, using only outputs already released, its own memory and its quantum references; it may entangle inputs across visits and across blocks, keep purifications, and act on a loan before returning it. It cannot read the device's inert initial purifier. The error is the largest trace distance, over all such observers, between the final state of the observer's systems with the device and with the ideal target: the strategy distance of~\cite{Chiribella08,GutoskiWatrous07}. For batch operation this is the half diamond distance.

\paragraph{The schedule.} Definition~\ref{to:def:schedule} below makes the schedule a first-class object: a finite ordered list of arrivals and releases of carriers the observer owns, each returned with the Hamiltonian it arrived with. Batch, real time, delayed release, interactive processes and loans are schedules. Schedules that create or retire carriers, or change Hamiltonians, need an explicit inventory term or a physical switch, and Section~\ref{to:sec:law} says exactly what.

\paragraph{The return rule, and why it has the form it has.} The controller returns \emph{approximately}: uniformly over observers, its final joint state with everything else is close in trace distance to its initial state tensored with the rest, and its actual mean energy and its actual entropy each change by $o(n)$. Trace-distance closeness alone would not do. As Faist, Oppenheim and Renner warned, a catalyst returned only approximately can embezzle work~\cite{FOR15}. Here is the loophole in its simplest form: a reference register in a broad, flat superposition of $n^2$ energy levels, shifted down by $n$ levels, is within trace distance about $\sqrt{2/n}$ of itself yet has handed over $n$ quanta of energy. Refunding energy alone is not enough either: a large mixed register can stay within trace distance $\eta$ of its initial state after every use while absorbing $n$ bits of entropy, with about $n/\eta$ qubits~\cite{Mausberg26}, and a controller built that way is a second bath. The lower bound of Theorem~\ref{to:thm:law} sees the controller only through $\Delta E_C-\kB T\ln2\,\Delta S_C$, so we require each term to be $o(n)$. The contract allows the controller's \emph{startup capital}, its initial free energy, to be as large as one likes and to grow with $n$; capital is invested and must come back, and it is not part of the rate. A controller that comes back only in state, as an embezzling state does, can lower the battery's work to the capacity or below it, but only by paying the difference from its own free energy, like a second battery; this is quantified in Appendix~\ref{app:batch}.

\paragraph{What the contract does not cover.} Resources that start correlated with the controller and must be returned in a prescribed joint state are not covered: such a borrowed catalyst would have to be modelled jointly with the controller, and the product initial state used throughout would fail. Promises of product inputs, restrictions on loans and frozen external fuel would be additional contract data; the observer class here has none of them. A task may, however, restrict its admissible inputs by a promise that is part of the task, and so narrow the observer class: the task of Theorem~\ref{thm:classes} is one, and its bounds are proved directly in Appendix~\ref{app:classes}. The lent schedule of Figure~\ref{fig:three}, arm~(ii) of Section~\ref{sec:protocol}, instead needs a restriction on loans, that the observer return the carrier unmodified; Section~\ref{sec:protocol} states it as a contract of its own.

\paragraph{Units.} Entropies are in bits and work in units of $\kB T\ln2$ in the main text; Appendices~\ref{app:law} and~\ref{app:module} use nats where the proofs are shorter.

Table~\ref{to:tab:contract} collects the contract, and Table~\ref{tab:exceptions} the principal departures of the thermodynamic results from it; Section~\ref{to:sec:beyond} fixes its own free operations, currencies and output inventory.
